
\documentclass[UTF8,12pt,a4paper]{ctexart}
\usepackage[margin=2.2cm]{geometry}
\usepackage{amsmath,amssymb,mathtools}
\usepackage{fontspec}
\usepackage{xcolor}
\usepackage{enumitem}
\usepackage{hyperref}
\usepackage{fancyvrb}
\usepackage{fvextra}
\usepackage{microtype}
\usepackage{titlesec}
\usepackage{fancyhdr}
\usepackage{setspace}

\hypersetup{
  colorlinks=true,
  linkcolor=blue!55!black,
  urlcolor=blue!55!black
}

\setstretch{1.18}
\setlist[itemize]{leftmargin=1.8em,itemsep=0.25em,topsep=0.3em}
\setlist[enumerate]{leftmargin=2.1em,itemsep=0.25em,topsep=0.3em}
\titleformat{\section}{\Large\bfseries}{\thesection}{0.8em}{}
\titleformat{\subsection}{\large\bfseries}{\thesubsection}{0.8em}{}
\titleformat{\subsubsection}{\normalsize\bfseries}{\thesubsubsection}{0.8em}{}
\pagestyle{fancy}
\fancyhf{}
\fancyfoot[C]{\thepage}
\renewcommand{\headrulewidth}{0pt}

\DefineVerbatimEnvironment{GuideVerbatim}{Verbatim}{
  breaklines=true,
  breakanywhere=true,
  fontsize=\small,
  frame=single,
  framesep=3mm,
  rulecolor=\color{gray!35},
  bgcolor=gray!7
}

\newcommand{\GuideInlineCode}[1]{\texttt{\detokenize{#1}}}

\begin{document}

\section*{七条路线总图：从问题到可改进的算法方案}

这份总图回答一个问题：比赛当天拿到一个混合整数带约束优化题，我们到底在脑子里怎么拆它？

七条路线不是七个互相孤立的模型。更像一条流水线：

\begin{GuideVerbatim}
业务问题
  -> 变量、目标、约束
  -> QUBO / Ising 能量函数
  -> 约束处理与可行性控制
  -> 退火 / QAOA / Hybrid 子问题求解
  -> 学习引导 warm-start / variable fixing / repair
  -> 解码、repair、benchmark、展示
\end{GuideVerbatim}

前六条是量子优化和量子启发优化的主线。第七条是学习引导层：它用 GNN、强化学习或轻量策略学习优化流程里的局部决策，但不替代可行性检查和最优性验证。

\subsection*{先抓住一个共同例子}

仓库里的 \texttt{data/sample\_problem.json} 可以当学习锚点。它的意思很简单：

\begin{itemize}
\item 有三个模型候选：\texttt{model\_a}、\texttt{model\_b}、\texttt{model\_c}。
\item 有两个增强选项：\texttt{boost\_x}、\texttt{boost\_y}。
\item 目标是收益最大。
\item 约束 1：三个模型里必须选且只选一个。
\item 约束 2：总资源不能超过预算。
\end{itemize}

这已经包含比赛里最常见的结构：选择变量、收益目标、互斥约束、预算约束、组合增益。

后面七条路线都可以围绕它理解：

\begin{itemize}
\item QUBO / Ising：把“选哪个更好”写成能量函数。
\item 约束处理：保证“只选一个”和“别超预算”不会被求解器破坏。
\item 退火：在能量地形里反复采样低能解。
\item QAOA：用参数化量子线路制造低能解的概率分布。
\item Constrained mixer：让量子线路尽量只在合法选择之间移动。
\item Hybrid：别把整个问题硬塞给量子模块，而是拆成经典部分和小 QUBO 子问题。
\item Learning-guided optimization：学习变量重要性、初始解、修复动作和参数调度，让前面几条路线更快找到好候选。
\end{itemize}

\subsection*{七条路线一句话}

路线 1：QUBO / Ising 建模

把离散优化问题改写成二次能量函数。低能量代表好解，变量只能是 0/1 或 spin。

路线 2：约束处理、编码、罚函数

让求解器别只追求低能量却给出非法答案。核心是 penalty、encoding、feasibility check、repair。

路线 3：Quantum Annealing / Simulated Annealing

把 QUBO 当成能量地形，通过退火过程找低能样本。它稳、好讲、适合做第一个能跑通的求解路线。

路线 4：QAOA / VQA

用量子门模型构造参数化线路，交替执行 cost layer 和 mixer layer，再用经典优化器调参数。

路线 5：Constrained Mixer / Warm-start / XY Mixer

不再只靠罚函数惩罚非法解，而是设计“只在可行空间里移动”的量子操作。它是理论亮点路线。

路线 6：Hybrid MILP / MIQP 分解

用经典优化处理大框架，把量子或量子启发算法用作小规模离散子问题求解器。它是最稳的比赛工程路线。

路线 7：Learning-Guided Optimization / 神经网络辅助混合优化

把 QUBO 或 MILP 表示成图，学习 warm-start、变量固定、branching、repair 和 QAOA/退火参数初值。它适合使用 4 张 64GB GPU，但它本质是启发式增强，不是最优性证明。

\subsection*{共同词表}

变量

你能控制的选择。例如选不选某个模型、某辆车走不走某条边、某台机组某时段开不开。

目标函数

你想优化的分数。例如最大收益、最小成本、最小延迟、最大服务质量。

约束

合法性规则。例如只能选一个、资源不能超、任务必须被完成、路径必须连续。

可行解

满足所有硬约束的解。比赛里要特别小心：很多量子或启发式方法会输出低能但不可行的 bitstring。

QUBO

Quadratic Unconstrained Binary Optimization。标准形式是：

\[
\begin{aligned}
\operatorname*{minimize} E(x) = c + \sum_{i} q_{i} x_{i} + \sum_{i<j} q_{ij} x_{i} x_{j}\\
x_{i} \in \{0, 1\}
\end{aligned}
\]

它叫 unconstrained，不是因为原问题没有约束，而是因为约束通常被塞进 penalty 里了。

Ising

物理里常用的 spin 形式：

\[
\begin{aligned}
\operatorname*{minimize} H(s) = C + \sum_{i} h_{i} s_{i} + \sum_{i<j} J_{ij} s_{i} s_{j}\\
s_{i} \in \{-1, +1\}
\end{aligned}
\]

QUBO 和 Ising 可以互转，常用关系是 \texttt{s = 2x - 1} 或 \texttt{x = (s + 1) / 2}。

Hamiltonian

在量子算法里，Hamiltonian 可以粗略理解成“能量函数的量子版本”。QAOA 的 cost Hamiltonian 就是在奖励低成本 bitstring。

Sampler

采样器。它不一定证明最优，只是反复吐出候选解和能量。退火器、QAOA 测量结果都可以看作 samples。

Ansatz

参数化量子线路的结构模板。QAOA 的 ansatz 由 cost layer 和 mixer layer 交替组成。

Mixer

让量子态在不同 bitstring 之间移动的操作。标准 X mixer 会独立翻转每个 bit；XY mixer 更像“把一个 1 移到另一个位置”，能保持固定选择数。

Relaxation

把离散变量放松成连续变量。例如把 \texttt{x in \{0,1\}} 放松成 \texttt{0 <= x <= 1}。它给出一个容易求的近似方向，但结果通常不是合法整数解。

Repair

把求解器输出的非法解修成合法解。比如预算超了，就删掉性价比最低的选择。

Warm-start

先给求解器一个较好的初始解、初始概率或初始参数。它不保证最优，但能减少盲搜。

Variable fixing

把高置信度变量先固定，只把不确定变量留给 QUBO、QAOA、退火或 hybrid 子问题。固定错了会伤害解，所以必须保留回退和 benchmark。

Learning-guided policy

学习一个局部策略，例如“哪个 bit 更可能为 1”“哪个变量适合分支”“下一步 repair 翻哪个变量”。它服务优化流程，而不是代替求解器。

\subsection*{一条路线的五个检查问题}

学任何一条路线时，都问这五个问题：

\begin{enumerate}
\item 搜索空间是什么？
\item 能量函数或目标函数怎么定义？
\item 算法允许怎么移动？
\item 如何保证或恢复可行性？
\item 比赛时我们能改哪些旋钮？
\end{enumerate}

这五个问题比背公式更重要。能回答它们，就能在赛题变形时改模型。

\subsection*{路线之间的关系}

路线 1 是语言层。

没有 QUBO / Ising，退火、QAOA 和很多 quantum-inspired 方法都缺输入。

路线 2 是安全层。

没有约束处理，算法可能很努力地找到一个低能量的非法答案。评委看到不可行解，不会因为你用了量子算法就加分。

路线 3 是第一个求解层。

模拟退火和量子退火都吃 QUBO/BQM。它们适合先做稳定 baseline。

路线 6 是工程扩展层。

一旦问题有连续变量、变量数太大、约束太多，就不要幻想直接全量 QUBO。先 hybrid，再把小离散块交给路线 3、4 或 5。

路线 4 是量子门模型展示层。

QAOA 很适合小规模 demo 和算法叙事，但要诚实面对 qubit 数、深度、shots、优化器不稳定等限制。

路线 5 是理论亮点层。

Constrained mixer 的价值是：有些约束不再靠巨大 penalty，而是由线路结构保持。这在 cardinality、one-hot、投资组合、分配问题里尤其有用。

路线 7 是学习加速层。

它从路线 1 的 QUBO 图、路线 2 的 violation、路线 3/4/5/6 的候选解里学习经验，再把 warm-start、变量固定、repair 或参数建议喂回求解流程。GPU 最适合用在这里。

\subsection*{比赛时怎么选路线}

如果你需要第一天就跑出稳定结果：

\begin{GuideVerbatim}
QUBO modeling + constraint checker + simulated annealing + repair
\end{GuideVerbatim}

如果赛题包含连续变量或大规模混合整数：

\begin{GuideVerbatim}
Hybrid relax-round-repair + small QUBO subproblem
\end{GuideVerbatim}

如果赛题有漂亮的 one-hot / exactly-k / portfolio 结构：

\begin{GuideVerbatim}
QAOA baseline + constrained XY mixer 对比
\end{GuideVerbatim}

如果评委更看重工程可靠性：

\begin{GuideVerbatim}
classical baseline + quantum-inspired solver + benchmark
\end{GuideVerbatim}

如果评委更看重量子算法亮点：

\begin{GuideVerbatim}
standard QAOA vs constrained QAOA vs SA
\end{GuideVerbatim}

如果比赛环境给了大显存 GPU：

\begin{GuideVerbatim}
QUBO graph dataset -> learning-guided policy -> warm-start / variable fixing -> SA / QAOA / Hybrid benchmark
\end{GuideVerbatim}

\subsection*{学习顺序}

推荐顺序是：

\begin{enumerate}
\item QUBO / Ising
\item 约束处理
\item 退火求解
\item Hybrid 分解
\item QAOA / VQA
\item Constrained mixer / warm-start
\item Learning-guided optimization
\end{enumerate}

这个顺序有点反直觉，因为 QAOA 和 GPU 听起来更“高级”。但真正比赛时，先懂建模和约束，后面所有量子算法和学习策略才有落脚点。

\subsection*{共读任务}

读完这份总图后，先别急着看论文。你可以先回答三个问题：

\begin{enumerate}
\item \texttt{sample\_problem.json} 里哪个约束最容易被算法破坏？
\item 如果 penalty 太小，会发生什么？如果太大，又会发生什么？
\item 如果模型数量从 3 个变成 300 个，你还会直接全量 QUBO 吗？
\item 如果神经网络很确信某些变量取值，你会直接相信它，还是把它当作 warm-start 和 variable fixing 建议？
\end{enumerate}

这几个问题答清楚，我们就可以进路线 1。

\subsection*{延伸阅读}

\begin{itemize}
\item \texttt{literature/notes/algorithm\_route\_classification.md}
\item \texttt{literature/notes/literature\_index.md}
\item \texttt{literature/notes/jerry\_reading\_plan.md}
\end{itemize}

\end{document}
